\subsection{Primitive elliptic bundles}
\label{sec:primitive-elliptic-bundles}

On a primitive elliptic bundle, the variation of the top-degree mass is a vertical Dirichlet integral. This gives the exact upper mass and a positive lower bound for separated potentials. The lower mass over all potentials remains a geometric question, since the collapse construction over a product of curves does not transfer directly to a base without curves.

Let $(B,\eta)$ be a connected compact K\"ahler surface and let $L\to B$ be a holomorphic line bundle such that
\[
 \int_Bc_1(L)\wedge\eta=0.
\]
Choose a smooth Hermitian metric on $L$ whose curvature is primitive with respect to $\eta$. Such a metric exists: for an arbitrary metric, the trace of the curvature with respect to $\eta$ has mean zero by the condition above, and the scalar Poisson equation gives a smooth change of metric that cancels this trace. Let $c$ be the resulting smooth closed curvature representative of $-c_1(L)$, so that $\eta\wedge c=0$. Write $\pi\colon L\to B$ for the bundle projection and $w$ for a nonzero vector in a fiber, with norm taken in the chosen metric. On the complement $L^\times$ of the zero section, put
\[
 t=\log|w|^2,\qquad \ddc t=\pi^*c.
\]
For $q>1$, form the quotient
\[
 X=L^\times/\langle w\mapsto qw\rangle,
 \qquad \pi\colon X\longrightarrow B,
 \qquad T=2\log q,
 \qquad A=\int_B\eta^2,
\]
where we use the same symbol $\pi$ for the induced projection. The action increases $t$ by $T$. It is free and properly discontinuous, and the compact annulus $1\leq|w|\leq q$ meets every orbit. Hence $X$ is a connected compact complex threefold; in this subsection it replaces the general manifold of the chapter. Over a trivializing open set $U\subset B$, a local holomorphic frame of $L$ identifies the fibration with $U\times(\C^*/q^{\mathbb Z})$.
\Needspace{5\baselineskip}
Since $\mathrm{d}t$ and $\dc t$ are invariant under the action, they descend to $X$, and we set
\[
 D=\mathrm{d}t\wedge\dc t,\qquad \omega=\pi^*\eta+D.
\]
The form $D$ has rank one; it is semipositive and positive on the vertical line, and its kernel projects isomorphically onto the tangent space of $B$. Hence $\omega$ is Hermitian.

\Needspace{17\baselineskip}
\begin{proposition}\label{prop:primitive-elliptic-metric}\index{Elliptic fibration!primitive bundles}
The metric $\omega$ is balanced and satisfies
\[
 \ddc\omega=-\pi^*(c^2)\geq0,\qquad
 \int_X\omega^3=3TA.
\]
For every smooth real $u$,
\begin{equation}\label{eq:primitive-elliptic-mass}
 \int_X(\omega+\ddc u)^3
 =3TA-3\int_X\mathrm{d}u\wedge\dc u\wedge\bigl(-\pi^*(c^2)\bigr).
\end{equation}
In particular $v_+(\omega)=3TA$. The canonical bundles satisfy $K_X\simeq\pi^*K_B$. If $c_1(L)\ne0$ in real cohomology, $X$ admits neither a pluriclosed metric nor a closed K\"ahler current. If $c_1(L)=0$ in real cohomology, the chosen $\omega$ is K\"ahler.
\end{proposition}
\begin{proof}
Write $a=\mathrm{d}t$ and $b=\dc t$. Then $\mathrm{d}a=0$, $\mathrm{d}b=\pi^*c$, $\dc a=-\pi^*c$ and $\dc b=0$. Consequently
\[
 D^2=0,\qquad \mathrm{d}D=-a\wedge\pi^*c,
 \qquad \ddc D=-\pi^*(c^2),
\]
\[
 \mathrm{d}(\omega^2)=-2\pi^*(\eta\wedge c)\wedge a=0,
\]
so $\omega$ is balanced. Since $c$ is primitive, its eigenvalues relative to $\eta$ are $\lambda$ and $-\lambda$ for a function $\lambda\geq0$, so $c^2=-\lambda^2\eta^2$. This proves the asserted sign of $\ddc\omega$.

To compute the volume, fix an elliptic fiber $E=\pi^{-1}(b_0)$, where $b_0\in B$. In a logarithmic fiber coordinate $x+\mathrm{i}y$, the real periods are $\log q$ and $2\pi$, while $t$ equals $2x$ up to an additive constant and $D|_E=\pi^{-1}\mathrm{d}x\wedge\mathrm{d}y$. Thus $\int_ED=T$. Since $\omega^3=3\pi^*(\eta^2)\wedge D$, the volume formula follows.

For the mass identity, expand the cube in \eqref{eq:primitive-elliptic-mass}. The terms with one and with three factors $\ddc u$ integrate to zero, because $\omega^2$ is closed and $\ddc u$ is exact. For the remaining term, put $S=\ddc\omega$. Integration by parts gives
\[
 \int_X(\ddc u)^2\wedge\omega
 =\int_Xu\,\ddc u\wedge S
 =-\int_X\mathrm{d}u\wedge\dc u\wedge S,
\]
where the last equality uses $\mathrm{d}S=0$. This proves the identity. Its last integral is nonnegative, and $u=0$ attains the upper bound $3TA$. Since smooth and bounded potentials have the same mass extrema, by \cref{thm:gaud-volume-mass-extrema}, this proves $v_+(\omega)=3TA$.

If $c_1(L)=0$ in real cohomology, then $c$ is exact and Stokes' theorem gives $\int_Bc^2=0$. The pointwise identity $c^2=-\lambda^2\eta^2$ then forces $c=0$, and $\omega$ is closed. Otherwise $S=-\pi^*(c^2)$ is nonzero, smooth and semipositive. For a pluriclosed metric $\gamma$, Stokes' theorem would give
\[
 0<\int_X\gamma\wedge S
 =\int_X\omega\wedge\ddc\gamma=0,
\]
a contradiction. Likewise, a closed positive current $R\geq\varepsilon\omega$ with $\varepsilon>0$ would give $0<\int_XR\wedge S=\int_X\omega\wedge\ddc R=0$. All these pairings are pairings of currents with smooth forms.

Finally, let $w$ be a local holomorphic fiber coordinate given by a frame of $L$, and let $\sigma$ be a local holomorphic two-form on $B$. The form $\pi^*\sigma\wedge\mathrm{d}w/w$ is invariant under the quotient action. Changing the frame of $L$ adds to $\mathrm{d}w/w$ a horizontal one-form, whose wedge product with $\pi^*\sigma$ vanishes. These local sections of $K_X$ therefore have the transition functions of $\pi^*K_B$, which proves $K_X\simeq\pi^*K_B$.
\end{proof}

\Needspace{14\baselineskip}
For separated potentials, positivity bounds both the amplitude and the downward slope of the vertical derivative. The sharp variance bound under these two constraints determines how much mass such a potential can lose. For $M,T>0$, set
\[
 F(M,T)=
 \begin{cases}
  T^2/12,&T\leq2M,\\
  M^2-4M^3/(3T),&T\geq2M.
 \end{cases}
\]

\begin{lemma}\label{lem:periodic-variance-bound}
If a smooth real $T$-periodic function $p$ has mean zero, $|p|<M$ and $p'>-1$, then
\[
 \frac1T\int_0^T p^2\,\mathrm{d}s<F(M,T).
\]
The supremum of the left side over this class is $F(M,T)$.
\end{lemma}
\begin{proof}
For $p=0$ the inequality holds because $F(M,T)>0$. Otherwise set $a=\max p>0$, $-b=\min p<0$ and $r=a+b$. Integrating $-p'<1$ over a forward interval $I$ from a maximum to the next minimum gives $r<|I|\leq T$; also $r<2M$, since $|p|<M$. The function $G(s)=(a-s)(s+b)$ is nonnegative on the range of $p$. Since $p$ has mean zero and $-p'<1$, integration over $I$ gives
\[
 abT-\int_0^T p^2\,\mathrm{d}s
 =\int_0^T G(p)\,\mathrm{d}s
 \geq-\int_I G(p)p'\,\mathrm{d}s
 =\frac{r^3}{6}.
\]
As $ab\leq r^2/4$, the average of $p^2$ is at most $r^2/4-r^3/(6T)$. This expression increases strictly for $0<r<T$, and its value at $r=\min\{2M,T\}$ is $F(M,T)$. Since $r<\min\{2M,T\}$, the strict inequality follows.

For sharpness, consider the following $T$-periodic profiles $p_0$. When $T\geq2M$, let $p_0$ be constant equal to $M$ and to $-M$ on two intervals of length $(T-2M)/2$ each, joined by a descent of slope $-1$ and length $2M$, and close the period by an upward jump. When $T\leq2M$, let $p_0(s)=T/2-s$ for $0\leq s<T$, with its upward jump at the end of the period. In both cases $p_0$ has mean zero, $|p_0|\leq M$, $p_0'\geq-1$ in the sense of distributions, and the average of $p_0^2$ is $F(M,T)$; both constructions apply when $T=2M$. Convolution with a nonnegative periodic approximate identity preserves the mean and the bound on the distributional derivative, and multiplication by $1-\varepsilon$ makes both inequalities strict. As the convolution scale and $\varepsilon$ tend to zero, the resulting smooth functions converge to $p_0$ in $L^2$, which proves that the supremum is $F(M,T)$.
\end{proof}

\Needspace{18\baselineskip}
\begin{proposition}\label{prop:primitive-elliptic-separated}\index{Monge--Ampere mass@Monge--Amp\`ere mass!separated potentials}
Suppose $c_1(L)\ne0$ in real cohomology, and put
\[
 \kappa=-\frac{\int_Bc^2}{A}>0.
\]
Every smooth potential $u=f(t)+\pi^*v$ with $f$ of period $T$, $v\in C^\infty(B,\R)$ and $\omega+\ddc u>0$ satisfies
\[
 \int_X(\omega+\ddc u)^3
 >3TA\bigl(1-\kappa F(\kappa^{-1/2},T)\bigr)>0.
\]
Let $\lambda\geq0$ be the nonnegative eigenvalue function of $c$ relative to $\eta$, so $c^2=-\lambda^2\eta^2$, and put $M_0=1/\max_B\lambda$. For the subclass $v=0$, the exact mass infimum is
\[
 3TA\bigl(1-\kappa F(M_0,T)\bigr)>0.
\]
It is approached by smooth potentials $f(t)$ with $\omega+\ddc f(t)>0$, and it is not attained by such potentials.
\end{proposition}
\begin{proof}
Put $p=f'$. The chain rule gives
\[
 \omega+\ddc u
 =\pi^*(\eta+pc+\ddc v)+(1+p')D.
\]
Positivity on the vertical line implies $p'>-1$. Over every point of $B$, the function $t$ takes every value modulo $T$, so restriction to $\ker\partial t$ shows that $\eta+pc+\ddc v$ is a K\"ahler form on $B$ for every value of $p$. Its square has integral
\[
 \int_B(\eta+pc+\ddc v)^2=A(1-\kappa p^2)>0,
\]
since primitivity removes the mixed term and Stokes' theorem removes the exact terms. Hence $|p|<\kappa^{-1/2}$.

Only the vertical derivative enters \eqref{eq:primitive-elliptic-mass}, because $-\pi^*(c^2)$ has top degree in the base directions. For a $T$-periodic function $F_0$, integration over a fiber gives $\int_E F_0(t)D=\int_0^T F_0(s)\,\mathrm{d}s$, and therefore
\[
 \int_X(\omega+\ddc u)^3
 =3TA\left(1-\frac\kappa T\int_0^T p^2\,\mathrm{d}s\right).
\]
Now apply \cref{lem:periodic-variance-bound} with $M=\kappa^{-1/2}$. The formula for $F$ gives $F(M,T)<M^2$, which proves positivity of the bound.

For $v=0$, diagonalizing $c$ with respect to $\eta$ shows that the horizontal positivity condition is exactly $|p|<M_0$; together with $p'>-1$ it also suffices for positivity of the full form. Moreover $\kappa$ is the $\eta^2$-average of $\lambda^2$, so $\kappa\leq M_0^{-2}$. \Cref{lem:periodic-variance-bound} with $M=M_0$ therefore proves the claimed positive lower bound and its nonattainment. The smooth approximating functions in its proof have mean zero and hence periodic primitives $f_\varepsilon$, and the potentials $f_\varepsilon(t)$ prove that the bound is sharp.
\end{proof}

The proof of the first bound does not show that it is sharp for separated potentials over an arbitrary base: positivity of the classes $[\eta]+p[c]$ alone does not provide one function $v$ that makes all the representatives $\eta+pc+\ddc v$ positive. The threefold of \cref{sec:vanishing-lower-mass} is the primitive elliptic bundle over $\mathbb P^1\times\mathbb P^1$ with $q=\mathrm e^{2\pi}$, so that $A=2$, $\kappa=1$ and $T=4\pi$. For it, \cref{prop:primitive-elliptic-metric,prop:primitive-elliptic-separated} reduce to \cref{prop:lower-mass-balanced-metric,prop:lower-mass-no-pluriclosed,prop:elliptic-radial-mass}, the bound is $8$ and it is sharp, and by \cref{thm:vanishing-lower-mass} its lower mass vanishes for every Hermitian form.

A base without curves rules out the averages of horizontal curves used in the known collapse example. The following construction uses the period theorem for K3 surfaces and the description of their K\"ahler cone \cite[Theorems 7.4.1 and 8.5.2 of the online version]{HuyK3}.

\begin{example}\label{ex:primitive-elliptic-k3}\index{Elliptic fibration!primitive K3 example}
There is a K3 surface $B$ with $\operatorname{Pic}(B)=\mathbb Z[L]$, $c_1(L)^2=-4$, and a K\"ahler form $\eta$ orthogonal to $c_1(L)$. It has no compact curves. The associated $X$ has trivial canonical bundle, a balanced metric of finite upper mass, and no pluriclosed metric or closed K\"ahler current. It has no divisors, and every compact irreducible curve is an elliptic fiber.
\end{example}
\begin{proof}
In a hyperbolic plane of the K3 lattice, choose a basis $e,f$ with $e^2=f^2=0$ and $e\cdot f=1$, and put $\ell=e-2f$. This vector is primitive and has square $-4$. Choose a period in $\ell^\perp$ outside the hyperplanes perpendicular to integral vectors not in $\mathbb Z\ell$. Such a period exists by Baire's theorem: the period domain in $\ell^\perp$, of signature $(3,18)$, is nonempty and each excluded hyperplane is proper. The period theorem and the description of integral $(1,1)$-classes give the asserted Picard group. An irreducible curve would have nonzero class, because its K\"ahler area is positive, and hence square $-4m^2<-2$ for some integer $m\neq0$, contradicting adjunction. Thus there are no curves.

The orthogonal complement of $\ell$ in $H^{1,1}(B,\R)$ has signature $(1,18)$. Choose a class of positive square there, with sign chosen so that it lies in the component of the positive cone containing the K\"ahler classes. Since there are no rational curves, the K\"ahler-cone theorem shows that this class is K\"ahler; choose $\eta$ in it. \Cref{prop:primitive-elliptic-metric} then gives the assertions about the metric and the canonical bundle.

If a divisor $Z\subset X$ dominated $B$, the push-forward of its integration current would be $\pi_*[Z]=m$ for a positive generic degree $m$. Stokes' theorem and the projection formula would then give
\[
 0=\int_X[Z]\wedge\ddc\omega
   =m\int_B(-c^2)>0.
\]
Thus no divisor dominates $B$. The image of any other irreducible divisor would be a curve or a point; the base has no curves, and a fiber has dimension one, so both alternatives are impossible. Likewise, the image of an irreducible compact curve is a point, and the curve is therefore an entire elliptic fiber.
\end{proof}

Put $V=3TA$. The root volumes of complementary tests can be computed explicitly for one radial family. For smooth real $T$-periodic functions $a,g$, consider
\[
 \Gamma=a(t)\eta^2+g(t)\eta\wedge D.
\]

\Needspace{15\baselineskip}
\begin{proposition}\label{prop:primitive-radial-root-volume}\index{Gauduchon volume!primitive radial tests}
The form $\Gamma$ is strictly positive and $\ddc$-closed exactly when $a$ is a positive constant and $g>0$. For such a form put $\bar g=T^{-1}\int_0^Tg(s)\,\mathrm ds$ and $d_\Gamma=\int_X\omega\wedge\Gamma$. Its positive square root and volume are
\[
 \gamma=\sqrt a\,\eta+\frac{g(t)}{2\sqrt a}D,
 \qquad
 \int_X\gamma^3=\frac32TA\sqrt a\,\bar g
       \leq\frac{d_\Gamma^{3/2}}{\sqrt V}.
\]
Equality holds exactly when $\bar g=2a$, without requiring $g$ to be constant. In particular every degree-one test in this family has root volume at most $V^{-1/2}$.
\end{proposition}
\begin{proof}
Since $\eta\wedge c=0$, the form $g\eta\wedge D$ is closed. The chain rule gives
\[
 \ddc(a\eta^2)=a''D\wedge\eta^2,
\]
the $a'c\wedge\eta^2$ term having excess base degree. Thus $\ddc\Gamma=0$ forces $a''=0$, and periodicity makes $a$ constant. In a positive coframe with $\eta=\sigma_1+\sigma_2$, $D=\sigma_3$, the complementary coefficient matrix of $\Gamma$ is $\operatorname{diag}(g,g,2a)$. Hence positivity is exactly $a,g>0$.
The displayed $\gamma$ is positive and squares to $\Gamma$. Multiplication and fiber integration give
\[
 d_\Gamma=TA(a+\bar g),\qquad
 \gamma^3=\tfrac32\sqrt a\,g\eta^2\wedge D.
\]
For fixed $d_\Gamma>0$ set $s=d_\Gamma/(TA)$. Then $0<a<s$, $\bar g=s-a$, and $\sqrt a(s-a)$ has its unique maximum at $a=s/3$. Substitution proves the bound and its equality condition.
\end{proof}

For the divisor-free threefold of \cref{ex:primitive-elliptic-k3}, \cref{thm:divisor-free-transverse-concentration} gives divergent transverse pairings against normalized Gauduchon tests. The radial family above has bounded root volume at fixed degree, even when $g$ concentrates arbitrarily strongly. Large transverse pairings do not by themselves give unbounded root volumes or a sequence of potentials with masses tending to zero.

\begin{problem}\label{prob:primitive-elliptic-lower-mass}
Determine the lower Monge--Amp\`ere mass (\cref{def:bmp}) of the primitive elliptic bundles other than the threefold of \cref{sec:vanishing-lower-mass}, for instance of their balanced metrics $\omega$ from \cref{prop:primitive-elliptic-metric}, in particular for the bundles of \cref{ex:primitive-elliptic-k3}. Does it vanish whenever $c_1(L)\neq0$ in real cohomology?
\end{problem}

The estimates for separated potentials do not answer this question; the known collapse over $\mathbb P^1\times\mathbb P^1$ uses potentials outside that class. On the K3 example every compact curve is vertical, so averaging over such curves cannot produce the horizontal directions used in \cref{lem:lower-mass-curve-average}. A proof of collapse would therefore need a different construction, while positive lower mass would require a bound for all admissible potentials or for all normalized Gauduchon tests.
